\printedpage{44}
% head: A. Grothendieck : Dérivateurs, Ch. XVII, Édition des Universités de Montpellier II et Paris VII
\[
\begin{tikzcd}
TC \arrow[r, hook] \arrow[d, leftrightarrow] & C \arrow[d, leftrightarrow] \\
F & TF \arrow[l, hook'] \ ,
\end{tikzcd}
\tag*{(5.2)}
\]
and $TC = C \cap W$, the fact that $(W, C, F)$ is a Quillen triple is corollary 4 p. 36 of XV: The conditions $TC = C \cap W$, $TF = F \cap W$, $W = TF \circ TW$ are equivalent, and equivalent to the fact that $(W, C, F)$ is a Quillen triple.

By Quillen's definition of triples that are \textit{closed,} by the properties
\[
\underbrace{W = TF \circ TC}_{\substack{\text{true for all}\\ \text{the Quillen}\\ \text{triples}}}, \qquad C = (TF)^*, \qquad \underbrace{F = (TC)_*}_{\substack{\text{true in}\\ \text{the present case}}}\ , % [sic] le triples -> les triples, agreement with "tous"
\]
it remains to express that $C = (TF)^*$, which is equivalent to $C$ stable under direct factors (cf. loc. cit. cor. 1, p. 31), q.e.d.

\textbf{Remark 5.4.} The statements 5.2 and 5.3, and the relation $T\tilde{C} \subset (TC)^\natural$ in proposition 5.1, do not call upon the whole hypothesis that $(\mathcal{M}, W, C)$ be a special category with cofibrations, but only that we have a category with cofibrations, with $C \mathrel{\leftarrow\!\!\!+\!\!\!\rightarrow} \underbrace{TF}_{\defeq C_*} \subset W$, % ?? uncertain: arrow symbol between C and TF -- double-headed arrow crossed by a small bar/plus sign
\msp{page 89}
i.e. a category with $W$, $C$ satisfying Cof$_0$, Cof$_1$, Cof$_2$ of page 39, plus $TF \defeq C_* \subset W$, plus the factorization axiom for $(C, TF)$ (which implies Cof$_3$).

\textbf{Corollary 5.5.} \textit{Let $(\mathcal{M}, W, C)$ be a special category with cofibrations, such that there exists a} small \textit{subset $TC_0 \subset TC = C \cap W$, with}
\[
TC \ \subset \ TC_0\cl\ ,
\tag*{(5.5.1)}
\]
\textit{or what amounts to the same thing, $(TC)\cl = (TC_0)\cl$, i.e.}
\[
\underbrace{F}_{\defeq (TC)_*} \ = \ (TC_0)_*\ .
\tag*{(5.5.2)}
\]
\textit{Then $(W, C, F)$ is a Quillen triple. This triple is closed if and only if $C$ is stable under direct factors.}

Indeed, the factorization theorem in XV §1 implies that under the hypothesis made on $TC_0 \subset TC$, the pair $(TC, F)$ satisfies the factorization condition. Indeed, by conditions b, b$'$ of definition 4.4, $TC$ has the desired stability properties (contains the isomorphisms, is stable under cobase change and under transfinite composition).

\msp{page 90}
\textbf{Definition 5.6.} The Quillen triples deduced as in 5.5 from a special category with cofibrations are called \textit{special.} Thus a Quillen triple $(W, C, F)$ in $\mathcal{M}$ is special if it satisfies conditions a) to d) in definition (4.4), concerning $(W, C)$ exclusively, plus the following two conditions.
% foot: Version 9 février 2019

\printedpage{45}
% head: A. Grothendieck : Dérivateurs, Ch. XVII, Édition des Universités de Montpellier II et Paris VII
\begin{itemize}[leftmargin=3em]
\item[\encircled{c$'$}] $\exists$ small subset $TC_0$ of $TC$, such that $TC \subset (TC_0)\cl$.
\item[\encircled{e}] $F = (TC)_*$.
\end{itemize}

\textbf{Corollary 5.7.} \textit{Let $(W, C, F)$ be a Quillen triple that is} closed \textit{in $\mathcal{M}$. For this triple to be special, it is necessary and sufficient that it satisfy the following conditions.}
\begin{itemize}[leftmargin=3em]
\item[\encircled{a}] \textit{$\mathcal{M}$ stable under small $\varinjlim$\,, and its elements are accessible.}
\item[\encoval{c, c$'$}] \textit{There exists a small subset $C_0$ of $C$, $TC_0$ of $TC$, such that we have}
\[
C \ = \ C_0\cl\ , \qquad TC \ = \ (TC_0)\cl\ .
\tag*{(5.7.1)}
\]
\end{itemize}
Indeed, all the other conditions a) to e), including c$'$) of 5.6, are automatically satisfied.

\msp{page 91}
\textbf{Proposition 5.8.} \textit{Let us consider a pair of adjoint functors}
\[
\begin{tikzcd}
\mathcal{M} \arrow[r, shift left, "v_!"] & \mathcal{M}' \arrow[l, shift left, "v^*"] \ ,
\end{tikzcd}
\tag*{(5.8.1)}
\]
\textit{where $\mathcal{M}$ is equipped with a special closed Quillen triple $(W, C, F = (TC)_*)$, with $TC = C \cap W$} (cf. definition 5.6). \textit{We suppose the following conditions satisfied} (compare prop. 4.6, relations (4.6.1), and (4.7.4)):
\[
\left\{
\begin{array}{ll}
\text{a)} & \mathcal{M}' \ \textit{stable under small}\ \varinjlim\ \textit{, its elements are accessible.}\\[1.5ex]
\text{b)} & v^* \ \textit{commutes with small}\ \varinjlim\ .\\[1.5ex]
\text{c)} & v^*(v_!(C)) \subset C,\ v^*(v_!(TC)) \subset TC.
\end{array}
\right.
\tag*{(5.8.2)}
\]
\textit{Let us define $(W', C', F')$ in $\mathcal{M}'$} (compare prop. 4.6, notably 4.6.2) \textit{by}
\[
\left\{
\begin{array}{lcl}
W' & = & (v^*)^{-1}(W)\\[1ex]
C' & = & v_!(C)\cl\\[1ex]
F' & = & (TC')_*\ , \qquad \textit{where}\ \ TC' \ \defeq \ C' \cap W'\ .
\end{array}
\right.
\tag*{(5.8.3)}
\]
\textit{Then $(W', C', F')$ is a closed special Quillen triple in $\mathcal{M}'$, and moreover we have}
\[
F' \ = \ (v^*)^{-1}(F)\ , \qquad TF' \ = \ (v^*)^{-1}(TF)\ ,
\tag*{(5.8.4)}
\]
\[
TC' \ = \ (v_!(TC))\cl\ .
\tag*{(5.8.5)}
\]

\msp{page 92}
\textsc{Proof}. By 4.6, we already know that $(\mathcal{M}', W', C')$ is a special category with cofibrations. (\textbf{N.B.} What is denoted $C'$ here, was denoted $\tilde{C}'$ in proposition 4.5. We must use the last assertion of 4.6, where we assume $C$ stable under direct factors, OK since $(W, C, F)$ assumed closed.) To prove that $(W', C', F' \defeq (TC')_*)$ is a Quillen triple
% foot: Version 9 février 2019
\printedpage{46}
% head: A. Grothendieck : Dérivateurs, Ch. XVII, Édition des Universités de Montpellier II et Paris VII
that is special, it remains, by 5.3 and definition 5.6, to prove that there exists a small subset $TC'_0$ of $TC'$, such that
\[
(\ TC'_0 \ \subset\ )\ TC' \ \subset \ (TC'_0)\cl\ .
\tag*{(5.8.6)}
\]
(This Quillen triple will moreover be closed, since $C'$ is Q-closed on the left by definition.) So let $TC_0$ be a small subset of $TC$ such that
\[
TC \ \subset \ (TC_0)\cl\ ,
\tag*{(5.8.7)}
\]
which exists by the hypothesis on $(W, C, F)$. We shall take
\[
TC'_0 \ \defeq \ v_!(TC_0)\ .
\tag*{(5.8.8)}
\]
We saw in 4.7 that we have
\[
v_!(TC) \ \subset \ TC'
\tag*{(5.8.9)}
\]
(follows from the hypothesis $v^*(v_!(TC)) \subset TC$ in (5.8.2. c)), a fortiori
\[
TC'_0 \ = \ v_!(TC_0) \ \subset \ TC'\ ,
\]
\msp{page 93}
it remains to prove
\[
TC' \ \overset{?}{\subset} \ (TC'_0)\cl\ ,
\tag*{(5.8.10)}
\]
which is equivalent to $(TC')\cl \subset (TC'_0)\cl$, or again to
\[
(TC'_0)_* \ \overset{?}{\subset} \ (TC')_* \ \defeq \ F'\ ,
\tag*{(5.8.10 bis)}
\]
whereas we have a priori the reverse inclusion (deduced from $TC'_0 \subset TC'$)
\[
F' \ \subset \ (TC'_0)_*\ .
\tag*{(5.8.11)}
\]
Now by adjunction, we have
\[
\begin{aligned}
(\ \underbrace{TC'_0}_{v_!(TC_0)}\ )_* \ &= \ (v_!(TC_0))_*\\
&= \ (v^*)^{-1}(\ \underbrace{(TC_0)_*}_{= (TC)_* \defeq F}\ )\\
&= \ (v^*)^{-1}(F)\ ,
\end{aligned}
\]
so (5.8.11) is written
\[
F' \ \subset \ (v^*)^{-1}(F)\ , \qquad \text{i.e.}\ \ v^*(F') \ \subset \ F
\tag*{(5.8.12)}
\]
(compare (4.7.6)), and we must show that this inclusion is an equality (on which I had gotten stuck p. 80, 81 \ldots), which is none other than the first of the two relations (5.8.4). But it is in the initial form (5.8.10) (transposed from (5.8.4)) that we are going to prove the thing, finally. Note that we have
% foot: Version 9 février 2019
\printedpage{47}
% head: A. Grothendieck : Dérivateurs, Ch. XVII, Édition des Universités de Montpellier II et Paris VII
\msp{page 94}
\[
(\underbrace{v_!(TC_0)}_{TC'_0})\cl \ = \ (v_!(TC))\cl\ ,
\]
since $TC$ is the closure of $TC_0$ under the operations of cobase change, of transfinite composition, and of direct factors, and $v_!$ commutes with these operations. So we must simply show
\[
TC' \ \subset \ (v_!(TC))\cl
\tag*{(5.8.12 bis)}
\]
(where there is no longer any question of $TC_0$). But the factorization theorem applies to the pair
\[
\Big( (v_!(TC))\cl,\ (v_!(TC))_* = (v^*)^{-1}(F) \Big)\ ,
\]
since $(v_!(TC))\cl = (TC'_0)\cl$, $TC'_0$ small. So if we have $f' \in TC' = C' \cap W'$, let us factor $f'$ as $p'i'$, with
\[
i' \ \in \ (v_!(TC))\cl\ , \qquad p' \ \in \ (v^*)^{-1}(F)\ .
\tag*{(5.8.13)}
\]
\[
\begin{tikzcd}[column sep=small, row sep=small]
X' \arrow[rr, "f' \in C' \cap W'"] \arrow[dr, "i'"'] & & Y' \arrow[dl, dotted, shift right, "q'"'] \\
& X'_1 \arrow[ur, shift right, "p'"'] &
\end{tikzcd}
\tag*{$(*)$}
\]
Note that we have
\msp{page 95}
\[
(v_!(TC))\cl \ \subset \ TC' \ (\ \defeq \ C' \cap W'\ )
\tag*{(5.8.14)}
\]
for $v_!(TC) \subset TC' \subset C'$, hence $(v_!(TC))\cl \subset \tilde{C}' = C'$, and it remains to prove that $v_!(TC)\cl \subset W'$, i.e. $v^*((v_!(TC))\cl) \subset W$. Now I say that we even have
\[
v^*((v_!(TC))\cl) \ \subset \ TC\ ,
\]
i.e.
\[
v_!(TC)\cl \ \subset \ (v^*)^{-1}(TC)\ .
\tag*{(5.8.15)}
\]
This follows from the fact that $TC$ is Q-closed, hence Q-saturated on the left, which implies, since $v^*$ commutes with small $\varinjlim$\,, that $(v^*)^{-1}(T)$ is so as well. On the other hand, it contains $v_!(TC)$, hence also the left Q-saturation of $v_!(TC)$. But the latter is \textit{equal} to the Q-closure $(v_!(TC))\cl$, by the factorization theorem, seeing that $TC'_0 \subset v_!(TC) \subset (TC'_0)\cl$, with $TC'_0$ small. % [sic] (v^*)^{-1}(T) -> (v^*)^{-1}(TC), the set in question is TC

Thus, (5.8.13) and (5.8.14) imply $i' \in TC' = C' \cap W'$, hence $i' \in W'$, whence $p' \in W'$, i.e. $v^*(p') \in W$, and since by hypothesis
% foot: Version 9 février 2019
\printedpage{48}
% head: A. Grothendieck : Dérivateurs, Ch. XVII, Édition des Universités de Montpellier II et Paris VII
\msp{page 96}
(5.8.13) we have $v^*(p') \in F$, we have $v^*(p') \in TF$, i.e.
\[
p' \ \in \ (v^*)^{-1}(TF) \ = \ TF' \ (\ \defeq \ C'\ )\ , % [sic] C' -> C'_*, cf. (5.8.16) TF' ( = C'_* )
\]
the equality $(v^*)^{-1}(TF) = TF'$, i.e.
\[
(v^*)^{-1}(C_*) \ = \ (v_!(C))_*\ ,
\]
resulting from the adjunction of $v_!$, $v^*$. Altogether, we therefore have
\[
p' \ \in \ TF' \ (\ = \ C'_*\ )\ ,
\tag*{(5.8.16)}
\]
and since $f' \in C'$, diagram $(*)$, page 114, shows us that there exists
\[
q' \ : \ Y' \longrightarrow X'_1\ ,
\]
with $p'q' = \mathrm{id}$, $q'f' = i'$. So $f'$ is a direct factor of $i' \in v_!(TC)\cl$, so we have $f' \in v_!(TC)\cl$, whence the relation
\[
TC' \ \subset \ v_!(TC)\cl\ ,
\]
which finishes proving that $(W', C', F')$ is a special and closed Quillen triple.

The two relations (5.8.4) were proved along the way (the second being essentially trivial, and the first representing the main difficulty in the proof). Remains to establish (5.8.5). But we already know (5.8.9) that
\[
v_!(TC) \ \subset \ TC'\ ,
\]
\msp{page 97}
whence
\[
v_!(TC)\cl \ \subset \ TC',
\]
since $TC'$ is Q-closed on the left (or because that was proved in (5.8.14)), we have already proved the inclusion in the reverse direction (5.8.12 bis), an assertion equivalent to (5.8.4) (and which represented the main difficulty).

\textbf{Remark 5.9.} If we do not assume that the structure on $\mathcal{M}$ is closed, then the proof given of 5.8 shows that we nevertheless find on $\mathcal{M}'$ a special Quillen structure, by defining $C'$ this time as in 4.6. But I doubt that we will ever have to work with special Quillen triples which are not also closed. In any case, if $W$ is stable under small filtering $\varinjlim$\footnote[28]{which implies stability under direct factors, and also implies the same property for $W' \defeq (v^*)^{-1}(W)$.} (and I doubt that one finds special Quillen triples outside this case), then if $(W, C, F)$ is a special Quillen triple, the same is true of $(W, \tilde{C} = C^\natural, F)$, and I presume that we have every interest in replacing the initial triple by this one. % footnote mark printed in the text as superscript (28)
% foot: Version 9 février 2019